% Values: PRECISE, FAITHFUL, ECONOMICAL
% Evidence planned:
% - PRECISE: each lexical class (reserved words, constants, comments,
%   identifiers, operators) defined by a stated rule; the three column tests
%   of the layout rule and every construct that sets the indentation given
%   exactly; one production for every bare Core construct of the book plan;
%   the grouping of operator chains stated as a rule, with the fixity of every
%   operator of the initial basis in a figure.
% - FAITHFUL: rulings A1, A2, A3, A4, A5, A6, A7, A8, A9, A10, A12, A13, A14,
%   A15, A16, A17, A26, A56 cited beside the text they shaped; every Core row
%   of the survey's "Syntactic Restrictions" table either enforced by the
%   grammar or lexis, or listed in Section 2.10, or left to the static
%   semantics because it needs types.
% - ECONOMICAL: bare grammar only (every derived form deferred to Appendix A);
%   tuple arity carried by the grammar rather than a restriction (A12);
%   restrictions already enforced by the grammar not repeated.

\chapter{Syntax of the Core}\label{ch:syntax-core}

\section{Reserved Words}

The following are the \emph{reserved words} of Standard Elm. They may not be
used as identifiers or as operators.
% A2, A16
\begin{quote}
\begin{tabular}{@{}*{7}{l}@{}}
\kw{as} & \kw{case} & \kw{else} & \kw{exposing} & \kw{if} & \kw{import} & \kw{in}\\
\kw{let} & \kw{module} & \kw{of} & \kw{port} & \kw{then} & \kw{type} & \kw{where}\\[0.5ex]
\kw{(}\ \kw{)} & \kw{[}\ \kw{]} & \kw{\kwlb}\ \kw{\kwrb} & \kw{,}\ \ \kw{.} & \kw{:}\ \ \kw{=} & \kw{|}\ \ \kw{->} & \kw{\textbackslash}\ \ \kw{\_}
\end{tabular}
\end{quote}
The words \kw{exposing}, \kw{import}, \kw{module}, \kw{port} and \kw{type} are
used by Modules (Chapter~\ref{ch:syntax-modules}). The word \kw{where} is used
by no phrase of Standard Elm: it belongs to the headers of effect modules,
which are kernel-only syntax (Section~\ref{sec:syntax-core-infix}). The words
\kw{alias}, \kw{infix}, \kw{left}, \kw{right}, \kw{non}, \kw{effect},
\kw{command} and \kw{subscription} are not reserved; \kw{alias} serves as a
terminal only directly after \kw{type} (Chapter~\ref{ch:syntax-modules}), and
elsewhere each of them is an identifier.

\section{Special Constants}\label{sec:syntax-core-constants}

Source text is a sequence of Unicode characters. A \emph{decimal digit} is one
of \kw{0}, \ldots, \kw{9}; a \emph{hexadecimal digit} is a decimal digit or one
of \kw{a}, \ldots, \kw{f}, \kw{A}, \ldots, \kw{F}.

An \emph{integer constant (in decimal notation)} is a non-empty sequence of
decimal digits that does not begin with \kw{0} unless it is \kw{0} itself. An
\emph{integer constant (in hexadecimal notation)} is \kw{0x} followed by a
non-empty sequence of hexadecimal digits. No integer constant carries a sign:
a negated integer is a derived form (Section~\ref{sec:syntax-core-lexical}).
An integer constant denotes a non-negative integer; its value as an
\kw{Int} is that integer normalised into the implementation's Int range by
$\norm$, the normalisation of Int arithmetic, so that every integer constant
has a value in $\IntRange$ (Appendix~\ref{ch:appendix-implementation}).
% A9, D8
Examples: \kw{0}\quad\kw{42}\quad\kw{0x1F}\quad\kw{0xffffffff}. Non-examples:
\kw{007}\quad\kw{0X1F}\quad\kw{1\_000}\quad\kw{0b101}.

\begin{figure}[tbp]
\centering
\begin{tabular}{@{}ll@{}}
\kw{\textbackslash n} & line feed (U+000A)\\
\kw{\textbackslash r} & carriage return (U+000D)\\
\kw{\textbackslash t} & horizontal tab (U+0009)\\
\kw{\textbackslash "} & \kw{"}\\
\kw{\textbackslash\textquotesingle} & \kw{\textquotesingle}\\
\kw{\textbackslash\textbackslash} & \kw{\textbackslash}\\
\kw{\textbackslash u\kwlb}$h$\kw{\kwrb} & the character with code point $h$,
  read in hexadecimal
\end{tabular}
\caption{Escape Sequences}\label{fig:core-escapes}
\end{figure}

A \emph{float constant} is an integer constant in decimal notation followed by
a fraction, by an exponent, or by a fraction and then an exponent. A
\emph{fraction} is \kw{.} followed by a non-empty sequence of decimal digits;
an \emph{exponent} is \kw{e} or \kw{E}, then optionally \kw{+} or \kw{-}, then
a non-empty sequence of decimal digits. The constant \kw{0} may be followed by
a fraction but not directly by an exponent. A float constant denotes the
rational number that it writes in decimal notation. Examples:
\kw{0.5}\quad\kw{3.14}\quad\kw{6.02e23}\quad\kw{1E-3}\quad\kw{0.0e1}.
Non-examples: \kw{.5}\quad\kw{1.}\quad\kw{1.e5}\quad\kw{0e5}.

A \emph{character constant} is a single quote \kw{\textquotesingle}, then
either one character other than \kw{\textquotesingle}, \kw{\textbackslash} and
line feed, or one escape sequence, then a single quote. It denotes the
character it contains, or the character its escape sequence denotes. Thus
\kw{\textquotesingle\textquotesingle} and
\kw{\textquotesingle ab\textquotesingle} are not constants.

A \emph{string constant (single-line)} is a double quote \kw{"}, then a
sequence of zero or more characters other than \kw{"}, \kw{\textbackslash} and
line feed, and escape sequences, then a double quote. Three double quotes
\kw{"""} always begin a multi-line string constant, so \kw{""} is the empty
string only when no third \kw{"} follows it. A \emph{string constant
(multi-line)} is \kw{"""}, then a sequence of characters other than
\kw{\textbackslash}, and escape sequences, that does not contain \kw{"""}, then
\kw{"""}. The first \kw{"""} after the opening one closes the constant; any
further \kw{"} immediately after it begins a new item.
% A15
Line feeds and carriage returns inside a multi-line string constant are
characters of the string, and so is a line feed directly after the opening
\kw{"""}. A string constant denotes the sequence of characters it contains,
each escape sequence standing for the character it denotes.

Each escape sequence begins with \kw{\textbackslash} and stands for one
character. The escape sequences are shown in Figure~\ref{fig:core-escapes},
subject to this restriction:
\begin{itemize}
\item\citelabel{res:core-escape-unicode} In \kw{\textbackslash u\kwlb}$h$\kw{\kwrb}, $h$ is a sequence of four, five or
  six hexadecimal digits whose value is a Unicode scalar value, that is, at most
  \kw{10FFFF} and outside the surrogate range \kw{D800} to \kw{DFFF}. No other
  sequence beginning with \kw{\textbackslash} is an escape sequence.
  % A8
\end{itemize}
A tab character may occur, unescaped, in a character or string constant. How a
character is represented as a value is an implementation parameter
(Appendix~\ref{ch:appendix-implementation}).
% A52, D8


We denote by \cls{SCon} the class of \emph{special constants}, that is, the
integer, float, character and string constants; we use \scon{} to range over
\cls{SCon}.

\section{Comments}

A \emph{line comment} is \kw{-{}-} and the characters after it up to the end of
the line. A \emph{block comment} is any character sequence within comment
brackets \kw{\kwlb-}~\kw{-\kwrb} in which comment brackets are properly
nested. No space is allowed between the two characters of a comment bracket.
A program in which a \kw{\kwlb-} is unmatched is not a phrase of the language.
A tab character may occur in a line comment but not in a block comment.

A comment bracket \kw{\kwlb-} directly followed by \kw{|} begins a \emph{doc
comment}, which extends, with comment brackets nested as in a block comment,
to its matching \kw{-\kwrb}. A doc comment is not a comment: it is an item of
the grammar of Modules, which admits it only after a module header and before
a top-level declaration (Chapter~\ref{ch:syntax-modules}). No doc comment
occurs in a Core phrase.
% A10

\section{Identifiers}

The classes of \emph{identifiers} for the Core are shown in
Figure~\ref{fig:core-identifiers}. We use \var, \tyvar{} and so on to range
over \cls{Var}, \cls{TyVar}, etc. For each class marked ``long'' there is a
class of \emph{long identifiers}: \cls{LongVar}, \cls{LongCon} and
\cls{LongTyCon}, over which \longvar, \longcon{} and \longtycon{} range. Their
syntax is
\begin{center}
\begin{grammar}
\gprod{\longvar}{\var}{identifier}
\galt{\modid\kw{.}\var}{qualified identifier}
\gprod{\longcon}{\con}{identifier}
\galt{\modid\kw{.}\con}{qualified identifier}
\gprod{\longtycon}{\tycon}{identifier}
\galt{\modid\kw{.}\tycon}{qualified identifier}
\end{grammar}
\end{center}
with no space or comment on either side of each \kw{.}. Throughout this
document the term ``identifier'', occurring without an adjective, refers to an
identifier that is not qualified.

A \emph{letter} is a character of one of the Unicode general categories Lu,
Ll, Lt, Lm and Lo; the version of the Unicode standard that assigns these
categories is an implementation parameter
(Appendix~\ref{ch:appendix-implementation}).
% A14
The identifiers are formed as this restriction states:
\begin{itemize}
\item\citelabel{res:core-identifier-characters} A \emph{lower identifier} is a letter of category Ll (which includes \kw{a},
  \ldots, \kw{z}) followed by a sequence of zero or more \emph{inner characters}:
  letters, decimal digits and underbars \kw{\_}. An \emph{upper identifier} is a
  letter of category Lu or Lt (which include \kw{A}, \ldots, \kw{Z}) followed by
  a sequence of zero or more inner characters. A letter of category Lm or Lo may
  continue an identifier but not begin one, and no identifier begins with
  \kw{\_}. In either case reserved words are excluded, and no component of a
  long identifier may be a reserved word: thus \kw{List.if} is not a long
  identifier.
  % A2
\end{itemize}
An \emph{operator} is a non-empty sequence of the \emph{symbols}
\begin{center}
\kw{+}\quad\kw{-}\quad\kw{*}\quad\kw{/}\quad\kw{=}\quad\kw{.}\quad\kw{<}\quad%
\kw{>}\quad\kw{:}\quad\kw{\&}\quad\kw{|}\quad\kw{\textasciicircum}\quad%
\kw{?}\quad\kw{\%}\quad\kw{!}
\end{center}
that is not a reserved word. This means that \kw{=}, \kw{:}, \kw{|} and
\kw{->} are not operators, but \kw{==}, \kw{::}, \kw{||} and \kw{|>} are.
Operators are never qualified.

A variable \var, a type variable \tyvar{} and a record label \lab{} are lower
identifiers. A constructor \con{} and a type constructor \tycon{} are upper
identifiers. A module identifier \modid{} is a non-empty sequence of upper
identifiers separated by \kw{.}, with no space or comment around each
\kw{.}. A type variable whose name begins with \kw{number}, \kw{comparable},
\kw{appendable} or \kw{compappend} is constrained to a class of types
(Section~\ref{sec:static-core-classes}).

\begin{figure}[tbp]
\centering
\begin{tabular}{@{}llll@{}}
\var & \cls{Var} & value variables & long\\
\con & \cls{Con} & value constructors & long\\
\tyvar & \cls{TyVar} & type variables & \\
\tycon & \cls{TyCon} & type constructors & long\\
\lab & \cls{Lab} & record labels & \\
\modid & \cls{ModId} & module identifiers & \\
\opv & \cls{Op} & operators &
\end{tabular}
\caption{Identifiers}\label{fig:core-identifiers}
\end{figure}

The classes \cls{Var}, \cls{TyVar} and \cls{Lab} share the lower identifiers,
and \cls{Con}, \cls{TyCon} and \cls{ModId} the upper identifiers. The syntax
class of an occurrence of an identifier in a Core phrase is determined thus:
\begin{enumerate}
\item An upper identifier immediately followed by \kw{.} is a component of a
  module identifier. The following rules assume that all occurrences of module
  identifiers have been removed.
\item Within a type expression, a lower identifier is a type variable and an
  upper identifier is a type constructor.
\item At the start of a component of a record expression, record update,
  record type or record pattern, and immediately after the \kw{.} of a field
  selection or field accessor (Appendix~\ref{ch:appendix-derived-forms}), a
  lower identifier is a record label. A record label that forms a component of
  a record pattern is also a variable, bound by that pattern.
\item Elsewhere, a lower identifier is a variable and an upper identifier is a
  constructor.
\end{enumerate}
For the remainder of this document we shall therefore assume that the classes
are all disjoint.

\section{Lexical Analysis}\label{sec:syntax-core-lexical}

Each item of lexical analysis is either a reserved word, a special constant, a
long identifier, an operator, or a doc comment. Comments and the
\emph{whitespace} characters (space, line feed and carriage return) separate
items and are otherwise ignored, except that the layout rule depends on the
column at which each item begins (Section~\ref{sec:syntax-core-layout}) and
that the constructs listed below depend on whether two items are separated. A
tab character may occur only within a line comment or within a character or
string constant. Where the remaining text begins with \kw{-{}-} or \kw{\kwlb-},
a comment or doc comment is taken. Otherwise, at each stage the longest next
item is taken.
% A3
Thus \kw{types} is an identifier, not \kw{type} followed by \kw{s}; within a
run of symbols \kw{-{}-} is part of an operator and begins no comment, so
\kw{a*-b} contains the single operator \kw{*-}, and \kw{2\textasciicircum-1}
the operator \kw{\textasciicircum-}. Items are also subject to these
restrictions:
\begin{itemize}
\item\citelabel{res:core-constant-separated} A special constant, and the reserved word \kw{\_}, may not be followed
  immediately by an inner character, and the digits of an integer constant may
  be followed immediately by \kw{.} only as part of a fraction; so \kw{12abc},
  \kw{1.x} and \kw{\_x} are not sequences of items.
\item\citelabel{res:core-adjacent-items} The following constructs require two items to be \emph{adjacent}, that is, to
  have no whitespace or comment between them:
  \begin{itemize}
  \item the parentheses of the unit \kw{()}, and the operator and parentheses of
    an operator value \kw{(}\opv\kw{)};
    % A1
  \item the \kw{.} of a long identifier and the items on either side of it;
  \item the \kw{.} of a field selection \atexp\kw{.}\lab{} and the items on
    either side of it, and the \kw{.} of a field accessor \kw{.}\lab{} and the
    label after it. A \kw{.} adjacent to a following label is the \kw{.} of a
    field selection only if it is adjacent to the item before it and that item
    ends an atomic expression that may precede a field selection
    (Section~\ref{sec:syntax-core-restrictions}); otherwise it begins a field
    accessor (Appendix~\ref{ch:appendix-derived-forms}). In particular a \kw{.}
    directly after a field accessor begins a field accessor;
    % A4
  \item the negation symbol and the atomic expression it negates.
  \end{itemize}
\end{itemize}
An item \kw{-} is a \emph{negation symbol} when it is adjacent to a following
atomic expression and either it begins an operand of an expression, or it
follows an operand from which it is separated by whitespace or a comment. A
negation symbol negates the atomic expression after it, and no more; in the
second case the negated expression is an argument of an application. Every
other \kw{-} is the operator \kw{-}. Thus \kw{f -x} applies \kw{f} to the
negation of \kw{x}, while \kw{a - b} and \kw{a-b} are subtractions; \kw{-f x}
applies the negation of \kw{f} to \kw{x}; and \kw{(-)} is the operator value
of subtraction. Negation is a derived form
(Appendix~\ref{ch:appendix-derived-forms}).

\section{Layout}\label{sec:syntax-core-layout}

The phrase structure of Elm depends on indentation. Source text is divided
into lines by line feeds. The \emph{column} of an item is one more than the
number of characters that precede it on its line, where a carriage return
counts as no character.
% A6
Each item is read under an \emph{indentation}, a natural number $i$. An item
whose column is $c$ is
\begin{itemize}
\item \emph{indented} if $c > i$ and $c > 1$;
\item \emph{aligned} if $c = i$;
\item \emph{at the margin} if $c = 1$.
\end{itemize}
These tests concern only the column at which an item begins: where the item
before it ended, for instance on a later line than it began, is irrelevant.
% A5
The \emph{layout rule} is this restriction:
\begin{itemize}
\item\citelabel{res:core-layout-rule} Every item of a phrase is indented, except the
  items that the constructs below require to be aligned or at the margin.
\end{itemize}
Only these constructs change the indentation; elsewhere an item is read under the
indentation of the phrase that contains it.
\begin{enumerate}
\item \emph{Top-level declarations} (Chapter~\ref{ch:syntax-modules}). The
  first item of each top-level declaration is at the margin, and its remaining
  items are read under the indentation $0$; so each of them lies in column $2$
  or beyond.
\item \emph{Let expressions.} Let $k$ be the column of the reserved word
  \kw{let}. The item after \kw{let} has a column $d$ with $d > k$ and $d > 1$,
  and $d$ is the indentation under which the definitions of the \kw{let} are
  read: the first item of each later definition, and of each annotation, is
  aligned, and every other item of the definitions is indented. The reserved
  word \kw{in} and the body after it are read under the indentation in force
  at \kw{let}.
\item \emph{Case expressions.} The items of \kw{case} \expm{} \kw{of} and the
  first item of the first match rule are read under the indentation in force
  at \kw{case}. The column $p$ of that first item is the indentation under
  which the match is read: the first item of each later match rule is
  aligned, and every other item of the match is indented.
\end{enumerate}
An annotation $\var\ \kw{:}\ \ty$ is followed directly by a definition of the
same \var, whose first item has the column of the annotation's first item: at
the margin at top level, aligned in a \kw{let}.
% A17
Each construct extends as far as possible: it takes in every following item
that passes the test that applies to it and continues the phrase, and it ends
before the first item that does not. Thus two definitions of one \kw{let}, or
two match rules of one \kw{case}, never begin on the same line, while
\kw{let x = 1 in x} is a phrase on one line.

\section{Infixed Operators}\label{sec:syntax-core-infix}

An operator \opv{} has a \emph{fixity}: a \emph{precedence}, which is one of
the digits \kw{0}, \ldots, \kw{9}, and an \emph{associativity}, which is
\emph{left}, \emph{right} or \emph{non}. Operators are declared, with their
fixities, only by kernel-only syntax. Kernel-only syntax, namely operator
declarations, effect module headers and references to \kw{Elm.Kernel}
modules, lies outside Standard Elm.
% A11, A16, A37
No program can therefore declare an operator, and the operators are those of
the initial basis, the same in every program. Their fixities are shown in
Figure~\ref{fig:core-fixities}; the operator table of the initial basis gives
the meaning of each, its type scheme in
Appendix~\ref{ch:appendix-static-basis} and its value in
Appendix~\ref{ch:appendix-dynamic-basis}. Operators of library modules beyond
the initial basis are outside this definition.
% Design decision 1 (audit X1, X5): operators are fixed by the operator table
% of the initial basis, not looked up in a context or environment.

\begin{figure}[tbp]
\centering
\begin{tabular}{@{}clll@{}}
Precedence & Left & Right & Non\\
\hline
9 & \kw{<{}<} & \kw{>{}>} & \\
8 & & \kw{\textasciicircum} & \\
7 & \kw{*}\ \ \kw{/}\ \ \kw{//} & & \\
6 & \kw{+}\ \ \kw{-} & & \\
5 & & \kw{++}\ \ \kw{::} & \\
4 & & & \kw{==}\ \ \kw{/=}\ \ \kw{<}\ \ \kw{>}\ \ \kw{<=}\ \ \kw{>=}\\
3 & & \kw{\&\&} & \\
2 & & \kw{||} & \\
0 & \kw{|>} & \kw{<|} &
\end{tabular}
\caption{Fixities of the Operators of the Initial Basis}\label{fig:core-fixities}
\end{figure}

An expression of the form $\expm_0\ \opv_1\ \expm_1\ \cdots\ \opv_n\ \expm_n$
($n \geq 1$), in which each $\opv_j$ is an operator and each operand $\expm_j$
is an application or an atomic expression, possibly negated, except that the
last operand may be any expression, is an \emph{operator chain}. Application
binds more tightly than any operator. A chain is grouped so that an operator
of higher precedence binds more tightly than one of lower precedence, subject
to this restriction:
\begin{itemize}
\item\citelabel{res:core-operator-chain} If $\opv_j$ and $\opv_l$ ($j < l$) of a chain have equal precedence and no operator of lower
  precedence stands between them, then either both associate to the left and
  the chain groups to the left, or both associate to the right and it groups to
  the right; otherwise the chain is not a phrase of the language.
\end{itemize}
Precisely, a
chain is the infixed application $\expm'\ \opv_j\ \expm''$ in which
$\opv_j$ has the least precedence in the chain and is, among the operators of
that precedence, the rightmost if they associate to the left and the leftmost
if they associate to the right; $\expm'$ and $\expm''$ are the chains, or
single operands, to the left and right of $\opv_j$, grouped by the same rule.
For example,
\begin{center}
\begin{tabular}{@{}lll@{}}
\kw{a |> f |> g} & parses as & \kw{(a |> f) |> g}\\
\kw{f <| g <| a} & parses as & \kw{f <| (g <| a)}\\
\kw{x ::\ ys ++ zs} & parses as & \kw{x ::\ (ys ++ zs)}\\
\kw{a < b == c} & is illegal & \\
\kw{f <{}< g >{}> h} & is illegal &
\end{tabular}
\end{center}

An operator value \kw{(}\opv\kw{)} is an atomic expression that denotes the
operator \opv{} as a value; it is the only way to use an operator other than
in a chain. There are no sections, and no operator is qualified. The infixed
application $\expm_1\ \opv\ \expm_2$ is a derived form
(Appendix~\ref{ch:appendix-derived-forms}); so are the infixed forms of
\kw{\&\&} and \kw{||}, which evaluate their second operand only when needed,
whereas the operator values \kw{(\&\&)} and \kw{(||)} are ordinary functions.
% A56
In the semantic rules every chain has been grouped and every infixed
application replaced by its equivalent form, since fixity affects only
parsing.

\section{Derived Forms}

Many syntactic forms of Elm have a meaning that is given by translation into a
smaller set of forms, called the \emph{bare language}. Among them are infixed
application, negation, conditional expressions, list expressions with
elements, list patterns with elements, field selection and field accessors,
lambda expressions and function bindings with parameters beyond the first,
and \kw{let} expressions with several definitions. These derived forms, and
their equivalent forms in the bare language, are given in
Appendix~\ref{ch:appendix-derived-forms}. Appendix~\ref{ch:appendix-full-grammar}
collects the bare and derived syntax into one grammar.

\section{Grammar}

The phrase classes for the Core are shown in
Figure~\ref{fig:core-phrase-classes}. We use the variable \atexp{} to range
over \cls{AtExp}, etc. The grammatical rules for the Core are shown in
Figures~\ref{fig:core-grammar-expressions}, \ref{fig:core-grammar-patterns}
and~\ref{fig:core-grammar-types}.

\begin{figure}[tbp]
\centering
\begin{tabular}{@{}ll@{}}
\cls{AtExp} & atomic expressions\\
\cls{ExpRow} & expression rows\\
\cls{Exp} & expressions\\
\cls{Match} & matches\\
\cls{MRule} & match rules\\[0.5ex]
\cls{Dec} & declarations\\
\cls{ValBind} & value bindings\\[0.5ex]
\cls{AtPat} & atomic patterns\\
\cls{PatRow} & pattern rows\\
\cls{Pat} & patterns\\[0.5ex]
\cls{Ty} & type expressions\\
\cls{TyRow} & type-expression rows
\end{tabular}
\caption{Core Phrase Classes}\label{fig:core-phrase-classes}
\end{figure}

The following conventions are adopted in presenting the grammatical rules,
and in their interpretation:
\begin{itemize}
\item The brackets $\opt{\ }$ enclose optional phrases.
\item Alternative forms for each phrase class are in order of decreasing
  precedence. L (resp.\ R) means left (resp.\ right) association.
\item The phrases $\kw{\textbackslash}\,\atpat\ \kw{->}\ \expm$,
  $\kw{case}\ \expm\ \kw{of}\ \match$ and $\kw{let}\ \dec\ \kw{in}\ \expm$
  extend as far right as the layout rule allows
  (Section~\ref{sec:syntax-core-layout}). Such a phrase is written in
  parentheses when it is the argument of an application, or an operand of an
  operator chain other than the last.
\item In a type construction $\longtycon\ \ty_1\ \cdots\ \ty_n$ with
  $n \geq 1$, each argument $\ty_j$ is a type variable, a type construction
  without arguments, a record type expression, a unit or tuple type
  expression, or a parenthesised type expression; thus \kw{Maybe List a}
  gives \kw{Maybe} two arguments. Type construction binds more tightly than
  \kw{->}.
\item In patterns, a constructed pattern binds more tightly than \kw{::},
  and \kw{::} more tightly than \kw{as}; so \kw{x ::\ xs as l} is a layered
  pattern whose pattern is \kw{x ::\ xs}. The \kw{::} of a pattern always
  denotes the list constructor of the initial basis. There is no negated
  constant in a pattern.
  % A7
\item The separator \kw{;} between match rules and the connective \kw{and}
  between value bindings are abstract syntax: neither is an item of source
  text, and \kw{and} is not a reserved word. In source text the match rules of
  a match, and the definitions of a \kw{let}, are separated by layout
  (Section~\ref{sec:syntax-core-layout}). Likewise the annotated value binding
  $\var\ \kw{:}\ \ty\ \kw{=}\ \expm$ is written in source text as the
  annotation $\var\ \kw{:}\ \ty$ followed by the definition
  $\var\ \kw{=}\ \expm$.
  % A17
\item In the bare language the declaration of a \kw{let} is a single value
  binding whose members form one group; a \kw{let} with several definitions
  is a derived form. A value binding $\atpat\ \kw{=}\ \expm$ in which \atpat{}
  is a variable is read by the first alternative for \valbind.
  % DECIDED: `x = e` matches both alternatives of valbind; it is read as a
  % variable binding, as the implementation tries a named definition first
  % (survey "Declarations": named definition then destructuring).
\item Tuples have exactly two or three components.
  % A12
\end{itemize}

\begin{figure}[tbp]
\centering
\begin{grammar}
\gprod{\atexp}{\scon}{special constant}
\galt{\longvar}{value variable}
\galt{\longcon}{value constructor}
\galt{\kw{(}\opv\kw{)}}{operator value}
\galt{\kw{()}}{unit}
\galt{\kw{(}\ \expm_1\ \kw{,}\ \expm_2\ \opt{\kw{,}\ \expm_3}\ \kw{)}}{tuple}
\galt{\kw{[}\ \kw{]}}{empty list}
\galt{\kw{\kwlb}\ \opt{\exprow}\ \kw{\kwrb}}{record}
\galt{\kw{\kwlb}\ \var\ \kw{|}\ \exprow\ \kw{\kwrb}}{record update}
\galt{\kw{let}\ \dec\ \kw{in}\ \expm}{local declaration}
\galt{\kw{(}\ \expm\ \kw{)}}{parentheses}
\noalign{\smallskip}
\gprod{\exprow}{\lab\ \kw{=}\ \expm\ \opt{\kw{,}\ \exprow}}{expression row}
\noalign{\smallskip}
\gprod{\expm}{\atexp}{atomic}
\galt{\expm\ \atexp}{application (L)}
\galt{\kw{\textbackslash}\,\atpat\ \kw{->}\ \expm}{function}
\galt{\kw{case}\ \expm\ \kw{of}\ \match}{case analysis}
\noalign{\smallskip}
\gprod{\match}{\mrule\ \opt{\kw{;}\ \match}}{match}
\noalign{\smallskip}
\gprod{\mrule}{\pat\ \kw{->}\ \expm}{match rule}
\noalign{\smallskip}
\gprod{\dec}{\valbind}{value declaration}
\noalign{\smallskip}
\gprod{\valbind}{\var\ \opt{\kw{:}\ \ty}\ \kw{=}\ \expm\ \opt{\kw{and}\ \valbind}}{variable binding}
\galt{\atpat\ \kw{=}\ \expm}{destructuring binding}
\end{grammar}
\caption{Grammar: Expressions, Matches, Declarations and Bindings}
\label{fig:core-grammar-expressions}
\end{figure}

\begin{figure}[tbp]
\centering
\begin{grammar}
\gprod{\atpat}{\kw{\_}}{wildcard}
\galt{\scon}{special constant}
\galt{\var}{variable}
\galt{\longcon}{constructor}
\galt{\kw{()}}{unit}
\galt{\kw{(}\ \pat_1\ \kw{,}\ \pat_2\ \opt{\kw{,}\ \pat_3}\ \kw{)}}{tuple}
\galt{\kw{[}\ \kw{]}}{empty list}
\galt{\kw{\kwlb}\ \opt{\patrow}\ \kw{\kwrb}}{record}
\galt{\kw{(}\ \pat\ \kw{)}}{parentheses}
\noalign{\smallskip}
\gprod{\patrow}{\lab\ \opt{\kw{,}\ \patrow}}{pattern row}
\noalign{\smallskip}
\gprod{\pat}{\atpat}{atomic}
\galt{\longcon\ \atpat_1\ \cdots\ \atpat_n}{constructed pattern ($n \geq 1$)}
\galt{\pat_1\ \kw{::}\ \pat_2}{list construction (R)}
\galt{\pat\ \kw{as}\ \var}{layered}
\end{grammar}
\caption{Grammar: Patterns}
\label{fig:core-grammar-patterns}
\end{figure}

\begin{figure}[tbp]
\centering
\begin{grammar}
\gprod{\ty}{\tyvar}{type variable}
\galt{\kw{\kwlb}\ \opt{\tyrow}\ \kw{\kwrb}}{record type expression}
\galt{\kw{\kwlb}\ \tyvar\ \kw{|}\ \tyrow\ \kw{\kwrb}}{extensible record type expression}
\galt{\kw{()}}{unit type expression}
\galt{\kw{(}\ \ty_1\ \kw{,}\ \ty_2\ \opt{\kw{,}\ \ty_3}\ \kw{)}}{tuple type expression}
\galt{\longtycon\ \ty_1\ \cdots\ \ty_n}{type construction ($n \geq 0$)}
\galt{\ty\ \kw{->}\ \ty'}{function type expression (R)}
\galt{\kw{(}\ \ty\ \kw{)}}{parentheses}
\noalign{\smallskip}
\gprod{\tyrow}{\lab\ \kw{:}\ \ty\ \opt{\kw{,}\ \tyrow}}{type-expression row}
\end{grammar}
\caption{Grammar: Type Expressions}
\label{fig:core-grammar-types}
\end{figure}

\section{Syntactic Restrictions}\label{sec:syntax-core-restrictions}

The restrictions below are conditions on a program as it is written, before
any derived form is replaced by its equivalent form; where a restriction
names a derived form, the form is that of
Appendix~\ref{ch:appendix-derived-forms}. A program that breaks one of them is
not a phrase of the language.
% DECIDED: restrictions are stated on source phrases, before translation of
% derived forms, because the translation erases distinctions they depend on
% (A26: `f = \n -> f n` and `f n = f n` have the same equivalent form) and
% can itself bind a label as a variable (the field accessor).

\emph{Records.}
\begin{itemize}
\item\citelabel{res:core-row-labels-distinct} No expression row, pattern row or type-expression row may contain the
  same \lab{} twice. This applies to the expression row of a record update.
\item\citelabel{res:core-record-update-form} In a record update $\kw{\kwlb}\ \var\ \kw{|}\ \exprow\ \kw{\kwrb}$ the
  item before \kw{|} is a variable that is not qualified, and at least one
  expression row follows \kw{|}; the grammar gives this shape, and an
  extensible record type expression likewise has at least one row.
\item\citelabel{res:core-field-selection} In a field selection $\atexp\kw{.}\lab$
  (Appendix~\ref{ch:appendix-derived-forms}), \atexp{} is a \longvar, a
  record, a record update, a field selection, or a phrase that begins with
  \kw{(}: a parenthesised expression, a tuple, the unit or an operator value.
  A field accessor is never followed by a field selection: in
  $\kw{.}\lab_1\kw{.}\lab_2$ the second \kw{.} begins a second field accessor,
  so the phrase is the application of $\kw{.}\lab_1$ to $\kw{.}\lab_2$.
  Likewise \kw{"s".x} is the application of a string constant to the field
  accessor \kw{.x}.
  % A4
\end{itemize}

\emph{Patterns and bound variables.} The variables \emph{bound} by a pattern
are the variables that occur in it, the labels of its pattern rows, and the
variables after its occurrences of \kw{as}. A \emph{binding site} is one of
the following: the parameters of a lambda expression, taken together; the
parameters of a function binding, taken together; the pattern of a match
rule; the pattern of a destructuring binding; and the definitions of one
\kw{let}, taken together, which bind the variables they define and the
variables their destructuring patterns bind. The \emph{scope} of a binding
site is the body of the lambda expression, the right-hand side of the function
binding, the expression of the match rule, or, for a destructuring binding
and for the definitions of a \kw{let}, the definitions of that \kw{let} and
its body.
\begin{itemize}
\item\citelabel{res:core-binding-site-distinct} No binding site may bind the same \var{} twice.
\item\citelabel{res:core-no-shadowing} No binding site may bind a \var{} that is already bound by a binding site
  in whose scope it lies. Constructors, type constructors and type variables
  are not subject to this restriction.
  % Audit X12: the top-level part of this restriction is stated in
  % Chapter syntax-modules only (S3).
\item\citelabel{res:core-no-float-pattern} No float constant may occur in a pattern.
\end{itemize}
The number of arguments in a constructed pattern is not restricted here: a
constructed pattern with too few or too many arguments for its constructor
does not elaborate (Chapter~\ref{ch:static-core}).
% Audit Y7: the arity restriction was removed; it needs the constructor's
% declaration (G6), and the static rules for constructed patterns enforce it.

\emph{Value bindings.}
\begin{itemize}
\item\citelabel{res:core-annotation-placement} In source text each annotation $\var\ \kw{:}\ \ty$ is followed directly
  by a definition of the same \var, laid out as
  Section~\ref{sec:syntax-core-layout} requires.
  % A17
\item\citelabel{res:core-let-recursion} In a \kw{let}, a definition \emph{refers to} a definition (possibly
  itself) if a variable that the second binds occurs in the first. Every
  definition that lies on a cycle of this relation is a function binding
  $\var\ \atpat_1\ \cdots\ \atpat_n\ \kw{=}\ \expm$ with $n \geq 1$; in
  particular, no destructuring binding and no definition
  $\var\ \kw{=}\ \expm$ without parameters lies on a cycle, even when \expm{}
  is a lambda expression. The corresponding restriction for top-level
  declarations is given in Chapter~\ref{ch:syntax-modules}.
  % A26
\end{itemize}

There is no limit on the size of a phrase: the number of parameters,
components, fields or constructors is unbounded, and the limits of an
implementation are listed in Appendix~\ref{ch:appendix-implementation}.
% A13, A34
